\chapter{Guide Dogs}

\section*{Definition and Overview}
Guide dogs are specially trained dogs that assist people who are blind or have low vision to move around safely and independently. They are trained to avoid obstacles, stop at kerbs and stairs, and travel safely in traffic, while their handlers make the decisions about where to go. Guide dogs are provided in Australia through organisations such as Guide Dogs Australia and Vision Australia, and they are matched carefully with their handlers. The guide dog and handler form a team that works, travels, and lives together. Guide dogs improve mobility, independence, confidence, and social connection for people with vision loss, and they are legally entitled to accompany their handlers into public places, including transport, shops, and restaurants.

\section*{Epidemiology}
There are estimated to be around 450,000 people in Australia with significant vision loss, of whom more than 50,000 are blind. Guide dogs assist a small but important proportion of these people, with organisations in Australia producing hundreds of guide dogs each year. A guide dog typically works for about 7--8 years before retirement, and handlers may use a guide dog for many years, sometimes with successive dogs. Demand for guide dogs generally exceeds supply, partly because of the cost and time involved in training, which takes around 1--2 years per dog.

\section*{Aetiology and Pathophysiology}
This chapter concerns an assistive service rather than a disease. The role and function of guide dogs include:

\begin{itemize}
    \item \textbf{Mobility assistance:} guiding the handler around obstacles and safely across roads
    \item \textbf{Orientation:} the handler directs the route while the dog ensures safe travel
    \item \textbf{Independence:} enabling travel alone, including to work, study, shops, and appointments
    \item \textbf{Safety:} stopping at kerbs, stairs, and hazards, and refusing to proceed into danger
    \item \textbf{Social connection:} reducing isolation and acting as a companion
    \item \textbf{Legal access:} under Australian law, guide dogs may accompany their handlers into public areas and transport
    \item \textbf{Matching:} dogs are matched to the handler's lifestyle, walking pace, and needs
\end{itemize}

\section*{Clinical Features}
People are considered for a guide dog when they have specific needs and circumstances.

\begin{itemize}
    \item Legal blindness or severe low vision that affects safe mobility
    \item The desire and ability to use a guide dog for daily travel
    \item Sufficient physical ability and orientation skills to direct the dog
    \item A lifestyle suited to dog ownership and care
    \item Situations where a guide dog would improve independence and quality of life
    \item Commitment to the training program, which includes a period of residential and in-home training
    \item Ongoing partnership, including regular assessments of the dog's health and work
\end{itemize}

\section*{Differential Diagnosis}
The decision to use a guide dog is individual and considers alternatives.

\begin{itemize}
    \item A white cane (long cane), which is the primary mobility aid for many people with vision loss
    \item Electronic travel aids and navigation apps
    \item Sighted guide assistance from other people
    \item Low vision rehabilitation, including orientation and mobility training
    \item Other assistance animals, such as therapy dogs (which have different roles)
    \item Alternative forms of transport and support services
\end{itemize}

\section*{When to Seek Medical Care}
People interested in a guide dog should discuss it with their eye specialist, orientation and mobility instructor, and the guide dog service. Medical review is needed for any change in vision, falls, or health problems affecting mobility.

\textbf{Call 000 (or local emergency services) for:}
\begin{itemize}
    \item Sudden loss of vision or a sudden change in vision, which needs urgent medical assessment
    \item A fall with injury or a suspected fracture
    \item Any medical emergency, including for the guide dog, where immediate help is needed
\end{itemize}

\section*{Diagnosis and Investigations}
Assessment for a guide dog is functional and practical.

\begin{itemize}
    \item \textbf{Vision assessment:} confirmation of the level of vision loss by an eye care professional
    \item \textbf{Orientation and mobility assessment:} evaluation of travel skills and readiness
    \item \textbf{Health assessment:} review of physical and mental health relevant to handling a dog
    \item \textbf{Lifestyle assessment:} the guide dog service considers the person's routine, home, and travel needs
    \item \textbf{Trial and training:} a structured training period in which the partnership develops
    \item \textbf{Ongoing support:} regular contact with the guide dog service, including veterinary care for the dog
\end{itemize}

\section*{Management}
Using a guide dog well is a partnership that involves training, care, and ongoing support.

\begin{itemize}
    \item \textbf{Training:} the handler learns to work with the dog, read its signals, and direct travel
    \item \textbf{Daily care:} feeding, grooming, exercise, and veterinary care for the dog
    \item \textbf{Regular practice:} maintaining skills with regular travel and route practice
    \item \textbf{Legal access:} knowing the rights and responsibilities of guide dog access
    \item \textbf{Health of the dog:} regular vet checks, and retirement of the dog at the appropriate time
    \item \textbf{Replacement:} a new dog is matched when the current dog retires
    \item \textbf{Community awareness:} educating others about not distracting or patting working guide dogs
    \item \textbf{Support services:} ongoing contact with the guide dog organisation
\end{itemize}

\section*{Complications}
\begin{itemize}
    \item The cost and time of training and caring for a dog
    \item Occasional incidents of public access being denied despite legal protections
    \item The dog's health or retirement disrupting established routines
    \item The emotional impact of retiring a working dog
    \item Situations where the dog is not suited to certain environments or travel
    \item Distraction of the dog by other people or animals
    \item The need to maintain the dog's training and handling skills
\end{itemize}

\section*{Prognosis and Outcomes}
Guide dogs have a profound effect on the lives of many people with vision loss. Studies show they increase independence, mobility, confidence, and social participation, and reduce reliance on others for travel. The guide dog partnership typically lasts around 7--8 years of working life, and organisations provide support for the dog's retirement and the transition to a successor dog. For people who are suitable, a guide dog is not just a mobility aid but a companion that enriches daily life.

\section*{Prevention and Self-Care}
\begin{itemize}
    \item Do not distract or pat a working guide dog while it is on duty
    \item Do not feed a guide dog without permission
    \item Allow guide dogs into all public places, including shops, transport, and restaurants
    \item Speak to the handler, not the dog, when communicating
    \item If interested in a guide dog, contact a guide dog organisation for an assessment
    \item If you have vision loss, have regular eye examinations and use low vision support services
    \item Keep up to date with vision rehabilitation, including orientation and mobility training
    \item Protect your eye health with regular check-ups and prompt treatment of eye conditions
\end{itemize}

\section*{Sources and Further Reading}
The following reputable sources provide further information for healthcare students and patients seeking reliable, current advice about guide dogs.

\begin{itemize}
    \item Guide Dogs Australia. \textit{Guide dog services and eligibility.} https://www.guidedogs.com.au/
    \item Vision Australia. \textit{Guide dogs and vision services.} https://www.visionaustralia.org/
    \item Seeing Eye Dogs Australia. \textit{Guide dog training and support.} https://sedac.org.au/
    \item Healthdirect Australia. \textit{Vision loss and support services.} https://www.healthdirect.gov.au/
    \item Australian Human Rights Commission. \textit{Guide dog access rights.} https://humanrights.gov.au/
    \item International Guide Dog Federation. \textit{Guide dog standards and information.} https://www.igdf.org.uk/
\end{itemize}
